% BLOCK15_STAGE2_MANAGED
% Lecture 15: Modern C++ after C++17 and next steps
\section{Лекция 15. Современный C++ после C++17}
\label{sec:lecture-15}

Основной стандарт курса --- C++17. К этому моменту у нас уже есть фундамент:
типы, lifetime, RAII, value semantics, smart pointers, move semantics,
templates, lambdas, STL, complexity и iterator invalidation.

Финальная лекция отвечает на другой вопрос: \textbf{что реально добавил C++20
поверх этой модели и зачем это понадобилось?}

\begin{importantbox}[Итог]
Новые стандарты не отменяют старые правила. Concept не исправляет плохой API,
range не отменяет lifetime, coroutine не означает отдельный thread, а module не
отменяет build graph. Современный C++ добавляет более точные инструменты поверх
прежних contracts.
\end{importantbox}

\subsection{Стандарт и реальная toolchain}

Нужно различать четыре вещи:
feature входит в стандарт; compiler реализует language feature; standard library реализует library component; build system умеет корректно собрать project с этой feature.

Поэтому \texttt{-std=c++20} --- необходимое, но не всегда достаточное условие.
Особенно это заметно на modules.

Поддержку проверяют по документации, feature-test macros и минимальному compile
probe в своей toolchain.

\subsection{Feature-test macros}

Пример:

\begin{cppcode}
#ifdef __cpp_concepts
static_assert(__cpp_concepts >= 201907L);
#endif
\end{cppcode}

Macro сообщает о feature support implementation. Конкретные значения следует
сверять с документацией/стандартными таблицами, а не запоминать навсегда.

\subsection{\texttt{constexpr}: compile-time возможность, не приказ}

\begin{cppcode}
constexpr int square(int value) {
    return value * value;
}

static_assert(square(5) == 25);

int input = read_value();
int runtime_result = square(input);
\end{cppcode}

Первый call обязан дать constant expression. Второй может быть обычным runtime
call.

\begin{warningbox}[constexpr не равно «всегда compile time»]
Keyword разрешает constant evaluation при выполнении правил языка. Он не
означает, что каждый вызов обязательно исчезнет из runtime.
\end{warningbox}

\subsection{Как развивался constexpr}

C++11 ввёл достаточно строгий вариант. C++14/17/20 постепенно разрешали всё
больше обычного imperative code и library operations в constant evaluation.

Это делает compile-time programming похожее на обычный C++, но I/O и
произвольные runtime resources не становятся compile-time операциями.

\subsection{\texttt{if constexpr}: напоминание из C++17}

\begin{cppcode}
template <typename T>
void print_value(const T& value) {
    if constexpr (std::is_integral_v<T>) {
        std::cout << "integer: " << value;
    } else {
        std::cout << value;
    }
}
\end{cppcode}

Discarded branch не требует той же well-formedness для текущего instantiation,
что обычная runtime ветка. Это важный мост от C++17 templates к C++20 Concepts.

\subsection{\texttt{consteval}: immediate function, C++20}

\begin{cppcode}
consteval int make_id(int group, int number) {
    return group * 1000 + number;
}

constexpr int id = make_id(12, 34);
\end{cppcode}

Если relevant call не может быть constant expression, program ill-formed:

\begin{cppcode}
int x = read_value();
// int id = make_id(1, x); // compile error
\end{cppcode}

\texttt{consteval} нужен, когда compile-time evaluation входит в
контракта, а не просто возможностью.

\subsection{\texttt{constinit}: другой contract}

\begin{cppcode}
constinit int global_counter = 0;
++global_counter; // допустимо
\end{cppcode}

\texttt{constinit} требует static initialization для объекта со static/thread
storage duration, но не делает object const.

\begin{center}
\begin{tabularx}{\textwidth}{L{0.20\textwidth}YY}
\toprule
Keyword & Смысл & Runtime mutation \\
\midrule
\texttt{constexpr} & constant-expression object/function capability & object нет \\
\texttt{consteval} & immediate function & относится к calls \\
\texttt{constinit} & static initialization requirement & может быть да \\
\bottomrule
\end{tabularx}
\end{center}

\subsection{Static initialization order}

\texttt{constinit} полезен, когда global/static object должен гарантированно
получить static initialization. Если initializer внезапно требует dynamic
initialization, получаем diagnostic.

Но keyword не делает global mutable state хорошим design автоматически.

\subsection{Почему появились Concepts}

Старый generic API:

\begin{cppcode}
template <typename T>
auto twice(const T& value) {
    return value + value;
}
\end{cppcode}

Из signature не видно requirement на \texttt{operator+}. Ошибка может проявиться
глубоко в instantiation.

До C++20 применяли documentation, \texttt{static\_assert}, type traits и
SFINAE/\texttt{enable\_if}. Concepts позволяют сделать requirements частью
interface.

\subsection{Named Concept}

\begin{cppcode}
template <typename T>
concept Numeric =
    std::integral<T> || std::floating_point<T>;
\end{cppcode}

Concept --- compile-time predicate над template arguments, участвующий в
constraint system.

\subsection{Constrained template}

\begin{cppcode}
template <Numeric T>
T twice(T value) {
    return value + value;
}
\end{cppcode}

Или:

\begin{cppcode}
template <typename T>
requires Numeric<T>
T twice(T value) {
    return value + value;
}
\end{cppcode}

Неподходящий type rejected на уровне constraints.

\subsection{Abbreviated function template}

\begin{cppcode}
auto twice(Numeric auto value) {
    return value + value;
}
\end{cppcode}

Это всё ещё compile-time generic function, а не dynamic typing.

\subsection{Requires-expression}

\begin{cppcode}
template <typename T>
concept HasSize = requires(const T& value) {
    { value.size() } -> std::convertible_to<std::size_t>;
};
\end{cppcode}

Expression не выполняется runtime: compiler проверяет, корректна ли требуемая
операция и удовлетворяет ли result constraint.

\subsection{Виды requirements}

Simple requirement:

\begin{cppcode}
requires(T value) {
    ++value;
    value++;
}
\end{cppcode}

Type requirement:

\begin{cppcode}
requires {
    typename T::value_type;
}
\end{cppcode}

Compound requirement:

\begin{cppcode}
requires(const T& value) {
    { value.size() } -> std::convertible_to<std::size_t>;
}
\end{cppcode}

Nested requirement:

\begin{cppcode}
requires {
    requires std::integral<T>;
    requires sizeof(T) <= sizeof(int);
}
\end{cppcode}

\subsection{Concepts не доказывают semantics}

Compiler может проверить наличие \texttt{a < b}, но не способен из одного
requires-expression доказать, что comparison задаёт корректный strict weak
ordering для всех values.

Semantic requirements STL из блока~\ref{sec:lecture-14} остаются частью
контракта programmer.

\subsection{Concepts и overload resolution}

\begin{cppcode}
template <typename T>
void describe(const T&) {
    std::cout << "generic";
}

template <std::integral T>
void describe(const T&) {
    std::cout << "integral";
}
\end{cppcode}

Для integral argument constrained overload может быть выбран как более
подходящий. Это compile-time selection, не virtual dispatch.

\subsection{Concept или static\_assert}

\texttt{static\_assert} хорош для invariants и локальных compile-time условий.

Concept хорош, когда requirement относится к template interface и должен
участвовать в overload resolution.

Они дополняют друг друга.

\subsection{Concept и runtime polymorphism}

Concept описывает compile-time structural requirements. Virtual base class из
лекции 7 --- runtime nominal polymorphism.

Если нужно хранить heterogeneous objects в runtime, Concept сам по себе не
заменяет virtual/type-erasure abstraction.

\subsection{Стоимость Concepts}

Сам constraint не добавляет runtime branch. Проверка происходит compile time.

Но templates/constraints имеют compile-time cost: instantiations, diagnostics,
build time и возможный рост code size для разных specializations.

\subsection{Ranges: следующий слой над STL}

C++17:

\begin{cppcode}
std::sort(values.begin(), values.end());
\end{cppcode}

C++20:

\begin{cppcode}
std::ranges::sort(values);
\end{cppcode}

Range interface уменьшает повторение begin/end и использует Concepts для
формализации требований.

\subsection{Ranges algorithms}

\begin{cppcode}
auto it = std::ranges::find(values, 42);
\end{cppcode}

Algorithms по-прежнему имеют complexity и iterator/range requirements. Новый
interface не отменяет фундамент блока 14.

\subsection{Projection}

\begin{cppcode}
std::ranges::sort(
    blocks,
    std::less{},
    &CourseBlock::number
);
\end{cppcode}

Projection отдельно говорит, какое поле считать sort key. Это часто проще
comparator lambda.

\subsection{Views}

\begin{cppcode}
auto selected =
    values
    | std::views::filter([](int x) {
          return x % 2 == 0;
      })
    | std::views::transform([](int x) {
          return x * x;
      });
\end{cppcode}

Views обычно описывают lazy processing, не создавая промежуточный container на
каждом шаге.

\begin{warningbox}[View не означает ownership]
View часто заимствует source. Source должен жить достаточно долго и не быть
invalidated. Ranges не отменяют lifetime.
\end{warningbox}

\subsection{Lazy evaluation}

Transformation обычно выполняется при iteration. Это может убрать лишние
copies/allocations и позволить early termination.

Но многоэтажный pipeline не всегда читабельнее простого loop. Modern syntax
нужно выбирать по ясности и требованиям.

\subsection{Borrowed ranges и dangling}

Ranges library умеет выражать некоторые случаи, где iterator в уничтоженный
temporary небезопасен, например через \texttt{std::ranges::dangling}.

Это полезная защита, но не полный lifetime verifier.

\subsection{\texttt{std::span}, C++20}

\texttt{span} --- non-owning view на contiguous sequence:

\begin{cppcode}
int sum(std::span<const int> values) {
    return std::accumulate(values.begin(), values.end(), 0);
}
\end{cppcode}

Caller может передать compatible array или vector без копирования элементов.

\begin{importantbox}[Span --- borrow]
\texttt{std::span} не владеет underlying storage и не продлевает lifetime.
Reallocation vector или уничтожение owner может сделать span dangling.
\end{importantbox}

\subsection{Static extent}

\begin{cppcode}
std::span<const int, 4> values;
\end{cppcode}

Размер может быть частью type, если fixed extent действительно часть contract.

\subsection{Designated initializers, C++20}

\begin{cppcode}
struct Config {
    int retries;
    bool verbose;
};

Config config{
    .retries = 3,
    .verbose = true
};
\end{cppcode}

Имена members повышают читаемость aggregate initialization.

В C++20 designators должны идти в declaration order. Это отличается от
некоторых C patterns.

\subsection{Defaulted \texttt{operator<=>}}

\begin{cppcode}
struct Version {
    int major;
    int minor;

    auto operator<=>(const Version&) const = default;
};
\end{cppcode}

Для подходящего value type compiler может синтезировать comparison по members.

Не каждому domain type нужен total ordering: semantics важнее удобного syntax.

\subsection{\texttt{std::erase\_if}}

C++17 erase-remove:

\begin{cppcode}
values.erase(
    std::remove_if(values.begin(), values.end(), predicate),
    values.end()
);
\end{cppcode}

C++20:

\begin{cppcode}
std::erase_if(values, predicate);
\end{cppcode}

Новый API короче, но понимание C++17 idiom объясняет устройство старого кода и
разделение algorithm/container mutation.

\subsection{\texttt{std::string::starts\_with}}

\begin{cppcode}
if (name.starts_with("lecture_")) {
    // ...
}
\end{cppcode}

Новые стандарты часто добавляют маленькие library utilities. Перед написанием
домашнего helper проверьте standard library нужной версии.

\subsection{\texttt{std::source\_location}}

C++20 даёт typed source-location API для logging/diagnostics:

\begin{cppcode}
void log(
    std::string_view text,
    std::source_location where =
        std::source_location::current());
\end{cppcode}

Это современная альтернатива многим macro-based wrappers, но logging policy и
стоимость всё равно проектируются отдельно.

\subsection{\texttt{std::jthread}}

\texttt{std::jthread} связывает thread lifetime с RAII: destructor join-ит
joinable thread. C++20 также добавляет cooperative stop tokens.

Но concurrency требует отдельного изучения memory model, data races, atomics и
synchronization. Одной секцией это не закрыть.

\subsection{Coroutines: что это}

C++20 добавляет:
\texttt{co\_await}; \texttt{co\_yield}; \texttt{co\_return}.

Coroutine --- function, execution которой может suspend/resume.

\begin{warningbox}[Coroutine не равно thread]
Coroutine сама по себе не обещает background execution, parallelism или async
I/O. Scheduler/executor/event loop --- отдельный слой.
\end{warningbox}

\subsection{Как compiler видит coroutine}

Compiler концептуально превращает function в state machine и создаёт coroutine
frame для состояния, которое должно пережить suspension.

Frame может содержать parameters, locals, bookkeeping и promise object.
Конкретный layout/allocation --- implementation detail semantics.

\subsection{Promise type}

Return type coroutine через coroutine protocol определяет
\texttt{promise\_type}. Promise сообщает compiler, как создать return object,
initial/final suspend, обработать yield/return и exception.

Поэтому low-level coroutine example быстро становится разбором protocol, а не
самой идеи suspend/resume.

\subsection{\texttt{co\_await}}

Awaiter protocol концептуально отвечает на три вопроса:
готов ли result, нужно ли suspend и что вернуть после resume.

В application code обычно работают с готовыми task/awaitable types framework,
а не пишут protocol вручную в каждом месте.

\subsection{\texttt{co\_yield}}

Generator-like coroutine выдаёт values по одному, сохраняя state между
возобновлениями. C++20 задаёт language mechanism, но не универсальный
\texttt{std::generator} container-like abstraction.

\subsection{Lifetime coroutine}

Coroutine frame --- resource. Нужно знать:
кто владеет handle/frame; когда вызывается destroy; какие references переживают suspension; живы ли owners borrowed objects.

Современный syntax снова приводит к ownership/lifetime.

\subsection{Modules: зачем}

Header model основан на textual inclusion и несёт repeated parsing, macro
leakage и сложные include dependencies.

C++20 modules вводят language-level import/export model:

\begin{cppcode}
export module math;

export int add(int a, int b) {
    return a + b;
}
\end{cppcode}

Consumer:

\begin{cppcode}
import math;
\end{cppcode}

\subsection{Module не равен precompiled header}

PCH оптимизирует traditional header model. Module имеет собственную language
semantics интерфейса/import.

Они не являются одним механизмом.

\subsection{Почему modules зависят от tooling}

Compiler/build system должен построить dependency graph module units и
intermediate module artifacts в корректном порядке.

Поэтому module workflow изучают вместе с конкретными версиями compiler и CMake.
В этом курсе modules остаются обзорной темой, а не фиктивной «одной командой».

\subsection{Headers продолжают жить}

Большой ecosystem остаётся header-based. Реальный project может смешивать
modules, classic libraries, header-only dependencies и C APIs.

Migration --- engineering decision.

\subsection{Allocators}

Containers имеют allocator parameter. Allocator отвечает за storage
acquisition/release policy, но не заменяет ownership model объектов.

\begin{cppcode}
std::vector<int, MyAllocator<int>> values;
\end{cppcode}

Полный allocator model включает traits и propagation rules; это отдельная advanced тема.

\subsection{\texttt{std::pmr}}

Polymorphic memory resources появились уже в C++17. Они позволяют выбирать
memory resource runtime и полезны для arena/monotonic allocation workloads.

Custom memory strategy имеет смысл после измерения реальной allocation problem,
а не «потому что быстрее».

\subsection{Compile-time cost}

Современные templates, Concepts и ranges могут увеличивать compiler work:
instantiations, constraint checking, сложные view types, diagnostics.

Это compile-time cost, не обязательно runtime cost.

\subsection{Runtime cost нельзя угадывать по modern syntax}

View может убрать промежуточный container, а может усложнить optimization.
\texttt{std::function} может ввести indirect call. Coroutine может иметь frame.

Performance проверяют profiling/benchmark/assembly, а не возрастом стандарта.

\subsection{Эксперимент 1: constexpr\_and\_concepts\_lab}

Каталог:
\filepath{examples/15\_modern\_cpp\_and\_next\_steps/01\_constexpr\_and\_concepts\_lab}.

Проверяется:
\texttt{constexpr} compile-time и runtime calls; \texttt{consteval}; mutable \texttt{constinit}; standard/custom Concepts; requires-expression; constrained overload; real compile rejection для неподходящего type; real compile rejection runtime argument у \texttt{consteval}.

\begin{shellcode}
cd examples/15_modern_cpp_and_next_steps/01_constexpr_and_concepts_lab
./run_checks.sh
\end{shellcode}

\subsection{Эксперимент 2: cpp20\_feature\_tour}

Каталог:
\filepath{examples/15\_modern\_cpp\_and\_next\_steps/02\_cpp20\_feature\_tour}.

Проверяется:
\texttt{std::ranges::sort} с projection; filter/transform views; \texttt{std::span<const int>}; designated initializers; defaulted \texttt{operator<=>}; \texttt{std::erase\_if}; \texttt{std::string::starts\_with}; real compile failure неправильного порядка designators.

\begin{shellcode}
cd examples/15_modern_cpp_and_next_steps/02_cpp20_feature_tour
./run_checks.sh
\end{shellcode}

\subsection{Почему coroutines/modules не основные experiments}

Concepts/ranges/constexpr хорошо проверяются обычным portable C++20 target.
Modules сильнее зависят от build workflow. Coroutine protocol требует
дополнительной library abstraction и легко уводит финальную лекцию в устройство
\texttt{promise\_type} вместо общей картины.

Поэтому их mechanics и границы разобраны честно, но без притворства, что тема
«освоена» одним игрушечным executable.

\subsection{Что появилось дальше: C++23}

После C++20 развитие продолжается. Например, C++23 включает:
\texttt{std::expected}; новые ranges facilities; \texttt{std::mdspan}; дальнейшие constexpr/library improvements.

Каждую такую feature нужно явно маркировать версией стандарта.

\subsection{\texttt{std::expected}: следующий error contract}

В блоке 10 были status/exceptions, в блоке 13 --- optional «есть/нет».

\texttt{expected<T,E>} выражает value \texttt{T} или error \texttt{E}. Это
следующий логичный шаг, когда caller обязан узнать причину failure.

\subsection{Как читать новую feature}

Рабочий алгоритм:
1) определить проблему старого подхода; 2) узнать standard/version; 3) прочитать semantics; 4) проверить ownership/lifetime; 5) понять compile/runtime cost; 6) сделать минимальный compile experiment; 7) только потом использовать в production.

\subsection{Что должно остаться после курса}

Нужно уметь объяснить:
declaration/definition и compile/link/runtime errors; object lifetime, ownership, pointer/reference/const contracts; RAII, copy/move и value categories; runtime и compile-time polymorphism; templates и Concepts; exception boundaries; lambdas/callables; STL container/iterator/algorithm contracts; complexity и invalidation; зачем ranges/Concepts появились после C++17.

\subsection{Дальше: concurrency}

Следующая самостоятельная область:
threads, mutexes, condition variables, atomics, happens-before, memory ordering,
lock-free structures, schedulers/executors.

Её нельзя безопасно изучить только через «запустили два thread».

\subsection{Дальше: performance engineering}

Profiling, hardware counters, cache behavior, branch prediction, vectorization,
assembly, allocation profiling и realistic benchmarks.

Сначала измерение, затем optimization.

\subsection{Дальше: library/API design}

ABI/API boundaries, pImpl, type erasure, customization points, error models,
ownership across library boundaries, allocator-aware design и versioning.

\subsection{Дальше: advanced templates}

Variadic templates, fold expressions, advanced traits, specialization,
constexpr algorithms и policy-based design.

Concepts помогают выразить requirements, но не отменяют эти механизмы.

\subsection{Дальше: tooling}

Warnings, tests, sanitizers, debugger, static analysis, formatter, CMake, CI и
profiling остаются частью профессионального C++, а не внешним дополнением к
языку.

\subsection{Что значит знать современный C++}

Не помнить все keywords, а уметь:
выбрать простейший корректный abstraction; понимать owner/lifetime; понимать value/reference semantics; читать diagnostics; знать complexity contract; различать compile-time и runtime mechanisms; проверять support стандарта; измерять performance claims.

\subsection{Совместимость проекта}

Если основной project обещает C++17, C++20 experiments отделяют явно:
отдельный CMake project/target; \texttt{cxx\_std\_20}; README с required standard; отсутствие C++20 API в C++17 targets.

\subsection{Финальный checklist новой feature}

Перед использованием спросите:
1) какую проблему решает; 2) в каком стандарте появилась; 3) поддерживает ли toolchain; 4) какие compile-time requirements; 5) какие lifetime/runtime consequences; 6) есть ли более простое решение; 7) улучшает ли новый вариант contract; 8) как его тестировать.

\begin{importantbox}[Финальный вывод курса]
C++ остаётся языком явных trade-offs. Новые стандарты дают больше способов
выразить намерение compile time и точнее ограничить generic code, но фундамент
остаётся тем же: lifetime, ownership, types, contracts, algorithms, complexity
и tooling. Если основы понятны, новые стандарты перестают быть списком
магических keywords.
\end{importantbox}

На этом базовая линия из 15 блоков завершена. Дальше выбирайте отдельную
ветку --- concurrency, performance, library design, networking, graphics,
embedded C++ или другую предметную область --- и строить её тем же способом:
theory, runnable experiments, diagnostics и проверяемые guarantees.
